\documentclass{beamer}
\usetheme{Madrid}
\usepackage[utf8]{inputenc}
\usepackage{amsmath}
\usepackage{amsfonts}
\usepackage{amssymb}
\usepackage{graphicx}
\usepackage{xcolor}
\usepackage{tikz}

\title{CAMS Exam Preparation}
\subtitle{Domain A: Risks \& Methods of Financial Crime (100 Study Slides)}
\author{Based on CAMS Candidate Handbook \& Candidate Feedback}
\date{}

\setbeamercovered{transparent}
\setbeamertemplate{navigation symbols}{}

\begin{document}
	
	% ============================================================
	% TITLE SLIDE
	% ============================================================
	\begin{frame}
		\titlepage
		\centering
		\textit{Domain A: 30\% of CAMS Exam | 100 Key Concepts \& Explanations}
	\end{frame}
	
	% ============================================================
	% SECTION A.1: DEFINITIONS OF AML, CFT, SANCTIONS, FRAUD, ABC, TAX EVASION
	% ============================================================
	
	\begin{frame}{A.1.1: Money Laundering vs. Terrorist Financing}
		\begin{block}{Key Distinction (Heavily Tested)}
			\begin{itemize}
				\item \textbf{Money Laundering}: Conceals \textbf{illicit origin} of proceeds; always involves dirty money from predicate crimes.
				\item \textbf{Terrorist Financing}: Funds may be \textbf{legitimate} (donations, salary); concealment objective is \textbf{intended use} for violence, not origin.
			\end{itemize}
		\end{block}
		\textcolor{red}{\textbf{CRITICAL:}} Standard AML controls focused on suspicious \textbf{origins} may fail to detect TF involving \textbf{clean-source funds}.
	\end{frame}
	
	\begin{frame}{A.1.2: FATF Recommendation 3 - Predicate Offenses}
		\begin{block}{Key Requirements}
			\begin{itemize}
				\item FATF recommends the \textbf{"all-crimes approach"} - applies to proceeds from \textbf{all serious criminal acts}.
				\item \textbf{Self-laundering IS criminalized} - predicate offender can also be charged with money laundering.
				\item 6AMLD harmonizes \textbf{22 predicate offenses} across EU.
			\end{itemize}
		\end{block}
		\textcolor{blue}{Example:} Drug trafficker who processes own proceeds through a car wash business $\rightarrow$ can be charged with BOTH drug trafficking AND money laundering.
	\end{frame}
	
	\begin{frame}{A.1.3: 6AMLD - Corporate Criminal Liability}
		\begin{block}{EU Sixth Anti-Money Laundering Directive (6AMLD)}
			\begin{itemize}
				\item \textbf{Harmonizes 22 predicate offenses} across all EU member states (binding).
				\item \textbf{Introduces criminal liability for legal persons (corporations)} - not just fines.
				\item Corporation can be prosecuted and subject to: fines, debarment, dissolution.
				\item Does NOT require conviction of individual directors first.
			\end{itemize}
		\end{block}
		\textcolor{red}{\textbf{TIP:}} Know the difference from 5AMLD (5AMLD added VASPs and lowered prepaid card thresholds).
	\end{frame}
	
	\begin{frame}{A.1.4: Consequences of Money Laundering}
		\begin{columns}
			\column{0.5\textwidth}
			\textbf{Economic Consequences}
			\begin{itemize}
				\item Distorts markets and asset prices
				\item Undermines legitimate private sector
				\item Loss of tax revenue
				\item Reputational damage
				\item Misallocation of resources
			\end{itemize}
			\column{0.5\textwidth}
			\textbf{Social Consequences}
			\begin{itemize}
				\item Funds criminal enterprises
				\item Corruption of public officials
				\item Erodes rule of law
				\item Fuels violence and terrorism
				\item Increases poverty
			\end{itemize}
		\end{columns}
	\end{frame}
	
	\begin{frame}{A.1.5: Consequences of Indiscriminate De-Risking}
		\begin{block}{De-Risking Definition}
			Termination of entire categories of customer relationships without individual risk assessment.
		\end{block}
		\textbf{Negative Consequences:}
		\begin{itemize}
			\item Legitimate businesses lose access to global financial system.
			\item Drives transactions into \textbf{unregulated, informal channels}.
			\item Harms financial inclusion and developing economies.
			\item Contradicts the mandatory \textbf{risk-based approach}.
		\end{itemize}
	\end{frame}
	
	\begin{frame}{A.1.6: Self-Laundering and Double Jeopardy (U.S. Law)}
		\begin{block}{Self-Laundering (18 U.S.C. \$1956)}
			\begin{itemize}
				\item U.S. law \textbf{DOES} criminalize self-laundering.
				\item Predicate offender CAN be charged with money laundering for processing own proceeds.
				\item \textbf{Double jeopardy argument fails} - each offense requires proof of different elements.
			\end{itemize}
		\end{block}
		\textcolor{blue}{Example:} Wire fraud + money laundering = TWO separate crimes.
	\end{frame}
	
	\begin{frame}{A.1.7: FATF Recommendation 5 - Terrorist Financing Offense}
		\begin{block}{TF Offense Requirements}
			\begin{itemize}
				\item Criminalize TF regardless of whether funds actually used for terrorist act.
				\item \textbf{No minimum monetary threshold} (even \$1 is enough).
				\item Collection of funds for terrorist purposes is sufficient.
				\item Standalone offense (not just covered under ML).
			\end{itemize}
		\end{block}
		\textcolor{red}{\textbf{Deficiency:}} \$10,000 threshold for TF violates Recommendation 5.
	\end{frame}
	
	% ============================================================
	% SECTION A.2: THREE-STAGE ANALYSIS
	% ============================================================
	
	\begin{frame}{A.2.1: The Three Stages of Money Laundering}
		\begin{block}{Three Stages}
			\begin{enumerate}
				\item \textbf{Placement}: Introducing illicit funds into the financial system.
				\item \textbf{Layering}: Conducting complex transactions to obscure the audit trail.
				\item \textbf{Integration}: Making laundered funds appear legitimate through investments.
			\end{enumerate}
		\end{block}
		\textcolor{blue}{Example:} Cash deposits $\rightarrow$ wire transfers $\rightarrow$ real estate purchase.
	\end{frame}
	
	\begin{frame}{A.2.2: Placement Techniques}
		\begin{block}{Common Placement Methods}
			\begin{itemize}
				\item Structured cash deposits (smurfing) below CTR thresholds
				\item Currency smuggling
				\item Purchasing assets with cash
				\item Using currency exchanges or MSBs
				\item Gambling at casinos
				\item Commingling with legitimate business revenues
			\end{itemize}
		\end{block}
		\textcolor{red}{Placement is the riskiest stage for criminals.}
	\end{frame}
	
	\begin{frame}{A.2.3: Structuring / Smurfing}
		\begin{block}{Structuring Definition}
			Deliberately breaking cash deposits into amounts below reporting thresholds to evade CTR requirements.
		\end{block}
		\textbf{Red flags:}
		\begin{itemize}
			\item Multiple deposits just below \$10,000 (\$9,500-\$9,900)
			\item Deposits from multiple individuals to same account
			\item Geographic dispersion of deposits
		\end{itemize}
		\textcolor{red}{Structuring is a federal crime independent of the underlying source of funds.}
	\end{frame}
	
	\begin{frame}{A.2.4: Layering Techniques}
		\begin{block}{Layering Methods}
			\begin{itemize}
				\item Wire transfers through multiple accounts and jurisdictions
				\item Use of shell companies and trusts
				\item Trade-based money laundering
				\item Cryptocurrency mixing and privacy coin conversion
				\item Use of professional intermediaries
			\end{itemize}
		\end{block}
		\textbf{Purpose:} Create distance between funds and criminal source.
	\end{frame}
	
	\begin{frame}{A.2.5: Integration Techniques}
		\begin{block}{Integration Methods}
			\begin{itemize}
				\item Purchase of real estate or luxury assets
				\item Investment in legitimate businesses
				\item False loan repayments
				\item Commingling with legitimate business revenues
				\item Sale of assets purchased with laundered funds
			\end{itemize}
		\end{block}
		\textbf{Purpose:} Funds re-enter economy as legitimate-appearing wealth.
	\end{frame}
	
	% ============================================================
	% SECTION A.3: TRADE-BASED MONEY LAUNDERING (TBML)
	% ============================================================
	
	\begin{frame}{A.3.1: Over-Invoicing (TBML Technique)}
		\begin{block}{Over-Invoicing Mechanism}
			\begin{itemize}
				\item Exporter inflates invoice price above market value.
				\item Importer pays inflated amount, transferring excess value to exporter.
			\end{itemize}
		\end{block}
		\textbf{Red flag:} Price significantly above market (e.g., \$1,200 vs \$180).
		\textcolor{blue}{Criminal purpose:} Move funds from importer to exporter disguised as trade.
	\end{frame}
	
	\begin{frame}{A.3.2: Under-Invoicing (TBML Technique)}
		\begin{block}{Under-Invoicing Mechanism}
			\begin{itemize}
				\item Exporter deflates invoice price below market value.
				\item Importer pays lower amount but sells goods at market price.
				\item Importer retains surplus as "clean" funds.
			\end{itemize}
		\end{block}
		\textbf{Red flag:} Price significantly below market.
		\textcolor{blue}{Criminal purpose:} Move value from exporter to importer.
	\end{frame}
	
	\begin{frame}{A.3.3: Phantom Shipping}
		\begin{block}{Phantom Shipping Mechanism}
			\begin{itemize}
				\item No actual goods are ever shipped.
				\item Complete fabricated documentation.
				\item Funds transferred based on false paperwork.
			\end{itemize}
		\end{block}
		\textbf{Detection indicators:}
		\begin{itemize}
			\item Empty containers discovered upon inspection.
			\item Shipper has no manufacturing capability.
			\item Altered bills of lading.
		\end{itemize}
	\end{frame}
	
	\begin{frame}{A.3.4: Free Trade Zone (FTZ) Vulnerabilities}
		\begin{block}{Why FTZs are High Risk}
			\begin{itemize}
				\item Relaxed regulatory oversight
				\item Minimal customs documentation inspections
				\item High volume of re-exported goods
				\item Opportunity for mislabeling cargo
				\item Phantom shipments
			\end{itemize}
		\end{block}
		\textbf{Red flags:} Uniform invoice values, empty containers, altered bills of lading.
	\end{frame}
	
	\begin{frame}{A.3.5: Additional TBML Techniques}
		\begin{block}{Multiple Invoicing and Short/Over-Shipping}
			\textbf{Multiple Invoicing:}
			\begin{itemize}
				\item Same goods invoiced multiple times through different banks
			\end{itemize}
			\textbf{Short-shipping / Over-shipping:}
			\begin{itemize}
				\item Discrepancy between invoice quantity and actual shipped goods
			\end{itemize}
			\textbf{Red flag:} Multiple L/C amendments without logical explanation.
		\end{block}
	\end{frame}
	
	% ============================================================
	% SECTION A.4: BLACK MARKET PESO EXCHANGE (BMPE)
	% ============================================================
	
	\begin{frame}{A.4.1: BMPE Mechanism - Complete Explanation}
		\begin{block}{Black Market Peso Exchange Steps}
			\begin{enumerate}
				\item Narco-dollars deposited in U.S. accounts (structured below \$10k).
				\item Colombian peso broker buys dollars at discount (20-30\% off).
				\item Broker uses dollars to pay U.S. invoices for Colombian merchants.
				\item Colombian merchant repays broker in pesos at favorable rate.
				\item \textbf{No cross-border currency movement} after initial deposit.
			\end{enumerate}
		\end{block}
		\textbf{Key insight:} Uses drug dollars to pay legitimate commercial invoices.
	\end{frame}
	
	\begin{frame}{A.4.2: BMPE Detection Red Flags}
		\begin{block}{Red Flags for BMPE}
			\begin{itemize}
				\item Structured cash deposits from multiple individuals.
				\item U.S. importer receiving payments from unrelated individuals.
				\item Foreign wholesaler receiving goods at below-market prices.
				\item No direct wire transfers between source and destination.
				\item Broker acting as intermediary with no apparent business purpose.
			\end{itemize}
		\end{block}
		\textbf{Detection challenge:} No cross-border movement; appears as domestic trade.
	\end{frame}
	
	% ============================================================
	% SECTION A.5: CORRESPONDENT BANKING
	% ============================================================
	
	\begin{frame}{A.5.1: Nested Correspondent Accounts}
		\begin{block}{Nested Correspondent Accounts - Critical Risk}
			When a respondent bank allows third parties to route transactions through the correspondent account \textbf{without disclosure}:
			\begin{itemize}
				\item Correspondent bank loses visibility into true originators/beneficiaries.
				\item Increased risk of processing for sanctioned entities.
				\item Violates core correspondent banking due diligence.
			\end{itemize}
		\end{block}
		\textbf{Required actions:} Immediate RFI, SAR filing, reassess relationship.
	\end{frame}
	
	\begin{frame}{A.5.2: Payable-Through Accounts (PTAs)}
		\begin{block}{PTA Obligations (U.S. Regulations)}
			\begin{itemize}
				\item Correspondent bank MUST apply CDD directly to underlying PTA accountholders.
				\item Must maintain records identifying each PTA accountholder.
				\item \textbf{Cannot rely solely on respondent bank's CDD.}
			\end{itemize}
		\end{block}
		\textbf{Red flag:} Multiple sub-\$10,000 wires to same shell company.
	\end{frame}
	
	\begin{frame}{A.5.3: Shell Bank Prohibition (Section 313)}
		\begin{block}{Section 313 - Shell Bank Definition}
			A "shell bank" has:
			\begin{itemize}
				\item NO physical presence in any country
				\item NOT affiliated with a regulated financial institution
			\end{itemize}
		\end{block}
		\textcolor{red}{\textbf{U.S. banks absolutely prohibited}} from maintaining correspondent accounts for foreign shell banks.
	\end{frame}
	
	\begin{frame}{A.5.4: Reverse Correspondent Banking Risks}
		\begin{block}{Reverse Correspondent Banking}
			Small bank providing account to larger international bank.
		\end{block}
		\textbf{Risks:}
		\begin{itemize}
			\item Small bank may be used as blind gateway.
			\item Small bank struggles with due diligence on complex structures.
			\item Creates moral hazard.
		\end{itemize}
	\end{frame}
	
	% ============================================================
	% SECTION A.6: VIRTUAL ASSETS AND VASPs
	% ============================================================
	
	\begin{frame}{A.6.1: FATF Travel Rule for VASPs}
		\begin{block}{Virtual Asset Travel Rule (Recommendation 16)}
			\textbf{Required information for transfers above USD/EUR 1,000:}
			\begin{itemize}
				\item Originator: name, account/wallet address
				\item Originator: address, ID number, or DOB (at least ONE)
				\item Beneficiary: name and account/wallet address
			\end{itemize}
		\end{block}
		\textcolor{red}{\textbf{Receiving VASP obligations:}} Must obtain info; if not provided $\rightarrow$ follow up, delay, apply EDD.
	\end{frame}
	
	\begin{frame}{A.6.2: DeFi and AML Responsibility}
		\begin{block}{FATF DeFi Guidance}
			If a person/entity maintains \textbf{control or sufficient influence} over a DeFi protocol:
			\begin{itemize}
				\item Through governance token majority (>50\%)
				\item Ability to upgrade contract code
				\item Ongoing management or control
			\end{itemize}
			\textcolor{red}{\textbf{They may be considered a VASP}} subject to AML obligations.
		\end{block}
	\end{frame}
	
	\begin{frame}{A.6.3: Privacy Coins (Monero, Zcash) Risk Features}
		\begin{block}{Anonymity-Enhanced Cryptocurrencies}
			\textbf{Monero (XMR):}
			\begin{itemize}
				\item Ring signatures - obscures which input is real
				\item Stealth addresses - one-time addresses per transaction
			\end{itemize}
			\textbf{Zcash (ZEC):}
			\begin{itemize}
				\item Zero-knowledge proofs (zk-SNARKs)
				\item Shielded transactions hide all details
			\end{itemize}
		\end{block}
	\end{frame}
	
	\begin{frame}{A.6.4: Crypto Mixers and Tumblers}
		\begin{block}{Mixing/Tumbling Process}
			\begin{enumerate}
				\item Customer sends crypto to mixing service.
				\item Mixer pools funds from many users.
				\item Mixer redistributes to new wallets controlled by customer.
				\item Original source becomes obscured.
			\end{enumerate}
		\end{block}
		\textbf{Red flags:} Use of mixer without legitimate explanation, refusal to explain source.
	\end{frame}
	
	\begin{frame}{A.6.5: NFT Wash Trading}
		\begin{block}{NFT Wash Trading Typology}
			\begin{enumerate}
				\item Self-controlled wallets buy/sell same NFTs among themselves.
				\item Prices artificially escalate through circular transactions.
				\item Creates fraudulent market record.
				\item Sells NFT to unwitting third party at inflated price.
				\item Converts to fiat as "capital gain."
			\end{enumerate}
		\end{block}
		\textbf{Exchange obligation:} File SAR based on wash trading pattern.
	\end{frame}
	
	\begin{frame}{A.6.6: Stablecoin Money Laundering Risks}
		\begin{block}{Stablecoin Vulnerabilities}
			\begin{itemize}
				\item Pegged to fiat currency - price stability facilitates value storage.
				\item Rapid cross-border transfers without banking intermediaries.
				\item Some issuers have weak or no CDD controls.
				\item Can be redeemed for fiat at any time.
			\end{itemize}
		\end{block}
		\textbf{Risk:} Cross-border value movement without traditional oversight.
	\end{frame}
	
	\begin{frame}{A.6.7: Custodial vs. Non-Custodial Wallets}
		\begin{block}{AML Obligations by Wallet Type}
			\textbf{Custodial wallets:}
			\begin{itemize}
				\item VASP has full CDD obligations
				\item VASP controls funds and can monitor transactions
			\end{itemize}
			\textbf{Non-custodial wallets:}
			\begin{itemize}
				\item VASP's CDD ability is limited
				\item Cannot verify off-platform transactions
				\item Higher inherent AML risk
			\end{itemize}
		\end{block}
	\end{frame}
	
	\begin{frame}{A.6.8: Cross-Chain Bridges and Layering}
		\begin{block}{Cross-Chain Bridge Exploitation}
			\begin{itemize}
				\item Criminal converts assets to token on privacy blockchain using bridge.
				\item Privacy blockchain obscures transaction trail.
				\item After mixing, funds bridged back as different asset.
				\item Transactional link to original source is severed.
			\end{itemize}
		\end{block}
		\textbf{Detection:} On-chain analysis can identify bridge usage patterns.
	\end{frame}
	
	% ============================================================
	% SECTION A.7: PEPs AND ENHANCED DUE DILIGENCE
	% ============================================================
	
	\begin{frame}{A.7.1: PEP Definition - Complete List}
		\begin{block}{Who Qualifies as PEP (FATF Definition)}
			\textbf{Individuals with prominent public functions:}
			\begin{itemize}
				\item Heads of state or government
				\item Senior politicians and cabinet ministers
				\item Senior judicial officials
				\item Senior military officials
				\item Senior executives of state-owned enterprises
				\item Leaders of political parties
			\end{itemize}
			\textbf{Family members and close associates} are also treated as PEPs.
		\end{block}
	\end{frame}
	
	\begin{frame}{A.7.2: PEP Family Members and Close Associates}
		\begin{block}{Family Member Definition}
			Treated as PEPs even without government positions:
			\begin{itemize}
				\item Spouse or civil partner
				\item Children (including adult children)
				\item Parents
				\item Siblings
			\end{itemize}
			\textbf{Close associates:}
			\begin{itemize}
				\item Close business or personal relationship to PEP
				\item Beneficial owners of entities for PEP's benefit
			\end{itemize}
		\end{block}
		\textcolor{red}{Example:} Adult daughter of foreign minister IS a PEP requiring EDD.
	\end{frame}
	
	\begin{frame}{A.7.3: PEP Lifecycle - Mid-Relationship Designation Change}
		\begin{block}{When a Customer Becomes a PEP}
			\begin{itemize}
				\item \textbf{Immediate reclassification to HIGH risk}
				\item Mandatory EDD applied within days
				\item Re-establish source of wealth AND source of funds
				\item Senior management approval to continue
				\item Back-dated gap period is control failure
			\end{itemize}
		\end{block}
		\textcolor{blue}{Example:} Customer elected to foreign cabinet $\rightarrow$ re-rate immediately.
	\end{frame}
	
	\begin{frame}{A.7.4: Source of Wealth vs. Source of Funds}
		\begin{block}{Critical Distinction (Heavily Tested)}
			\begin{itemize}
				\item \textbf{Source of Funds (SoF)}: Specific transaction being deposited (sale contract, inheritance).
				\item \textbf{Source of Wealth (SoW)}: How PEP accumulated overall net worth (career, business, inheritance).
			\end{itemize}
		\end{block}
		\textcolor{red}{\textbf{EDD for PEPs requires BOTH.}} Civil servant with \$4.5M apartment: SoF insufficient without SoW.
	\end{frame}
	
	\begin{frame}{A.7.5: Former PEP - Grace Period and Risk Assessment}
		\begin{block}{PEP Status After Retirement}
			\begin{itemize}
				\item Maintain for \textbf{"reasonable period"} (minimum 12 months)
				\item \textbf{Risk-based decision} considering:
				\begin{itemize}
					\item Country corruption risk
					\item Continued access to influence
					\item Time elapsed since leaving office
					\item Account activity pattern
				\end{itemize}
				\item 5 years clean activity MAY justify lowering risk
				\item \textbf{Decision must be documented}
			\end{itemize}
		\end{block}
	\end{frame}
	
	\begin{frame}{A.7.6: Domestic vs. Foreign PEPs - Differences}
		\begin{block}{Classification Differences}
			\textbf{Foreign PEP:}
			\begin{itemize}
				\item Always requires EDD and senior approval
				\item Higher inherent risk
				\item Must establish SoW and SoF
			\end{itemize}
			\textbf{Domestic PEP:}
			\begin{itemize}
				\item May require EDD depending on jurisdiction
				\item Potentially less intensive
				\item Still requires enhanced scrutiny
			\end{itemize}
		\end{block}
	\end{frame}
	
	\begin{frame}{A.7.7: EDD - Required Actions for PEPs}
		\begin{block}{Enhanced Due Diligence Requirements}
			\begin{itemize}
				\item Establish SoW and SoF through verifiable documentation
				\item Obtain senior management approval
				\item Increased transaction monitoring frequency
				\item Review all significant transactions
				\item Identify UBO of any legal entity or trust
				\item Document all findings
				\item Reassess risk annually or upon triggers
			\end{itemize}
		\end{block}
	\end{frame}
	
	% ============================================================
	% SECTION A.8: LIFE INSURANCE AML
	% ============================================================
	
	\begin{frame}{A.8.1: Life Insurance CDD Triggers}
		\begin{block}{When CDD Applies to Life Insurance}
			\begin{itemize}
				\item Annual premium exceeds USD/EUR 10,000
				\item Single premium exceeds USD/EUR 10,000
				\item Any suspicious transaction
			\end{itemize}
			\textbf{Who must be identified:}
			\begin{itemize}
				\item Policyholder
				\item Beneficiary (once designated)
				\item For legal entities: beneficial owners
			\end{itemize}
		\end{block}
		\textcolor{red}{Beneficiary CDD applies even if designated after policy issuance.}
	\end{frame}
	
	\begin{frame}{A.8.2: Life Insurance ML Exploitation Methods}
		\begin{block}{How Insurance Features Are Exploited}
			\textbf{Policy loan as layering:}
			\begin{itemize}
				\item Purchase policy with criminal proceeds
				\item Take policy loan for 90\% of cash value
				\item Loan proceeds are "legitimate" check
			\end{itemize}
			\textbf{Cash surrender extraction:}
			\begin{itemize}
				\item After brief holding period
				\item Surrender policy, receive check minus penalty
				\item Penalty is "laundering fee"
			\end{itemize}
		\end{block}
	\end{frame}
	
	\begin{frame}{A.8.3: Life Insurance AML Red Flags}
		\begin{block}{Life Insurance ML Indicators}
			\begin{itemize}
				\item Premium from offshore shell company
				\item Premium from secrecy haven jurisdiction
				\item Early surrender with proceeds to different account
				\item Policyholder is foreign PEP; beneficiary is shell company
				\item Large premium relative to legitimate income
				\item Premiums from multiple unrelated third parties
			\end{itemize}
		\end{block}
	\end{frame}
	
	% ============================================================
	% SECTION A.9: MSBs, AGENTS, AND PREPAID ACCESS
	% ============================================================
	
	\begin{frame}{A.9.1: MSB Registration Triggers (U.S. FinCEN)}
		\begin{block}{Who Must Register as an MSB}
			\begin{itemize}
				\item Check casher: cashing > \$1,000/person/day
				\item Money transmitter: any volume
				\item Currency dealer or exchanger
				\item Issuer or seller of prepaid access
				\item Traveler's check issuer
			\end{itemize}
		\end{block}
		\textcolor{red}{Grocery store cashing checks > \$1,000/day without registration = BSA violation.}
	\end{frame}
	
	\begin{frame}{A.9.2: MSB Principal Responsibility for Agents}
		\begin{block}{Agent Network Compliance}
			\begin{itemize}
				\item MSB principal is \textbf{fully responsible} for agent compliance
				\item Must monitor agent transaction activity
				\item Unexplained volume spikes require investigation
				\item May need to file SAR on agent location
				\item Terminate agent if compliance cannot be ensured
			\end{itemize}
		\end{block}
		\textbf{Red flag:} Agent processes 5,000\% above historical average.
	\end{frame}
	
	\begin{frame}{A.9.3: Unregistered MSB Operation}
		\begin{block}{Violations of Unregistered Check Casher}
			Convenience store cashing checks, "No ID Required":
			\begin{itemize}
				\item Registration violation: > \$1,000/person/day
				\item CTR violation: no filings over \$10,000
				\item SAR violation: no filings for suspicious patterns
				\item "No ID Required" violates CIP/CDD
			\end{itemize}
		\end{block}
	\end{frame}
	
	\begin{frame}{A.9.4: Prepaid Access Programs - AML Obligations}
		\begin{block}{Prepaid Access Requirements}
			\begin{itemize}
				\item Issuer is considered an MSB
				\item Must register with FinCEN
				\item Must implement risk-based AML program
				\item CDD on cardholders exceeding thresholds
				\item Monitor reload activity for structuring
				\item File CTRs and SARs as required
			\end{itemize}
		\end{block}
		\textbf{Not exempt:} Prepaid cards are monetary instruments under BSA.
	\end{frame}
	
	\begin{frame}{A.9.5: Prepaid Cards - Structuring and Border Evasion}
		\begin{block}{Prepaid Card Abuse Example}
			Traveler with 50 prepaid cards each loaded with \$9,500 (total \$475,000):
			\begin{itemize}
				\item \textbf{Structuring violation}: Each card just below \$10,000
				\item \textbf{Currency reporting evasion}: Cards are monetary instruments
				\item \textbf{Layering}: Purchased at multiple locations
				\item \textbf{Destination risk}: Flying to heightened scrutiny jurisdiction
			\end{itemize}
		\end{block}
	\end{frame}
	
	% ============================================================
	% SECTION A.10: HUMAN TRAFFICKING FINANCIAL INDICATORS
	% ============================================================
	
	\begin{frame}{A.10.1: Human Trafficking - Controller Account Pattern}
		\begin{block}{Trafficking Red Flags (Heavily Tested)}
			\begin{itemize}
				\item Multiple individuals sending small payments to single controller account
				\item Controller pays utilities for multiple residential addresses
				\item Weekly cash withdrawal of 85-95\% of balance
				\item No employment contracts or payroll records
				\item Victims share same surname as controller
			\end{itemize}
		\end{block}
		\textbf{Trafficking vs. Smuggling:} Trafficking = ongoing exploitation; Smuggling = one-time fee.
	\end{frame}
	
	\begin{frame}{A.10.2: Human Trafficking - Massage Parlor / Nail Salon Network}
		\begin{block}{Service Business Trafficking Indicators}
			\begin{itemize}
				\item Multiple small electronic payments (\$50-\$150) from dozens of individuals daily
				\item Funds transferred to central LLC holding commercial properties
				\item Properties operate as massage parlors or nail salons
				\item LLC pays mortgages and utilities
				\item Operators have no verifiable income other than deposits
			\end{itemize}
		\end{block}
		\textbf{Pattern:} Human trafficking for commercial sexual exploitation.
	\end{frame}
	
	% ============================================================
	% SECTION A.11: CASINO AND GAMBLING AML
	% ============================================================
	
	\begin{frame}{A.11.1: Casino AML Obligations (U.S. BSA)}
		\begin{block}{Casino Reporting Requirements}
			Casinos with revenue > \$1 million must:
			\begin{itemize}
				\item File CTRs for cash > \$10,000 in a gaming day
				\item File SARs for suspicious patterns > \$5,000
				\item Conduct CDD at \$3,000 chip purchase (FATF Rec. 12)
				\item Maintain AML compliance program
				\item Verify customer identity at threshold
			\end{itemize}
		\end{block}
	\end{frame}
	
	\begin{frame}{A.11.2: Chip Washing Scheme (Casino)}
		\begin{block}{Chip Washing Pattern}
			\begin{itemize}
				\item Large cash chip purchase (e.g., \$300,000) - triggers CTR
				\item \textbf{Minimal gaming} (e.g., 45 minutes, \$1,500 wagered)
				\item Cash out remaining chips for casino check
				\item Accepts small "gaming loss" as laundering fee
			\end{itemize}
		\end{block}
		\textbf{Obligations:} File CTR AND file SAR.
	\end{frame}
	
	\begin{frame}{A.11.3: Online Gaming Money Laundering}
		\begin{block}{In-Game Currency Laundering Stages}
			\begin{enumerate}
				\item \textbf{Placement}: Purchase in-game currency with criminal cash via prepaid cards
				\item \textbf{Layering}: Distribute across multiple gamer accounts
				\item \textbf{Integration}: Sell currency peer-to-peer for PayPal or bank transfer
			\end{enumerate}
		\end{block}
		\textbf{Difficulty:} Small transactions fall below reporting thresholds.
	\end{frame}
	
	% ============================================================
	% SECTION A.12: REAL ESTATE MONEY LAUNDERING
	% ============================================================
	
	\begin{frame}{A.12.1: FinCEN Geographic Targeting Orders (GTOs)}
		\begin{block}{GTO Requirements - Real Estate}
			\begin{itemize}
				\item Title companies must identify natural persons behind legal entities
				\item All-cash purchases above threshold (typically \$300,000)
				\item Targets specific metro areas (Miami, LA, NYC, etc.)
				\item \textbf{Temporary orders}, renewed periodically
			\end{itemize}
		\end{block}
		\textbf{Red flags:} Recently formed LLC, IOLTA wire, registered agent refuses UBO disclosure.
	\end{frame}
	
	\begin{frame}{A.12.2: Real Estate as Integration Vehicle}
		\begin{block}{Example Integration Scheme}
			\begin{enumerate}
				\item \textbf{Placement}: Structured cash deposits below \$10,000
				\item \textbf{Layering}: Funds into anonymous LLC
				\item \textbf{Layering/Integration}: LLC purchases property with cash
				\item \textbf{Integration}: Property generates rental income reported on taxes
				\item \textbf{Integration}: Property sold, proceeds deposited as legitimate
			\end{enumerate}
		\end{block}
		\textbf{Red flags:} Rapid appreciation, original anonymous purchase.
	\end{frame}
	
	% ============================================================
	% SECTION A.13: GATEKEEPERS (LAWYERS, ACCOUNTANTS, TCSPs)
	% ============================================================
	
	\begin{frame}{A.13.1: IOLTA Account Vulnerabilities}
		\begin{block}{Why Lawyer Trust Accounts Are Dangerous}
			\begin{itemize}
				\item Banks rely on attorney's professional status
				\item Do NOT look through pooled accounts
				\item Illicit funds enter credentialed by attorney status
				\item Multiple client funds commingled
				\item Rapid disbursement possible
			\end{itemize}
		\end{block}
		\textbf{Red flags:} Wire from shell company, disproportionate retainer, no face-to-face contact.
	\end{frame}
	
	\begin{frame}{A.13.2: Trust and Company Service Provider (TCSP) Risks}
		\begin{block}{TCSP AML Obligations (FATF Rec. 23)}
			TCSPs must:
			\begin{itemize}
				\item Conduct CDD when forming legal persons
				\item Identify beneficial owners
				\item Maintain records for 5+ years
				\item File SARs when suspicious
			\end{itemize}
		\end{block}
		\textbf{Red flags:} Creates shell companies, nominee director, no physical office, refuses UBO disclosure.
	\end{frame}
	
	% ============================================================
	% SECTION A.14: TRUSTS AND BEARER SHARES
	% ============================================================
	
	\begin{frame}{A.14.1: Trust Beneficial Ownership - Who to Identify}
		\begin{block}{Trust CDD Requirements (FATF Rec. 25)}
			Must identify and verify:
			\begin{itemize}
				\item \textbf{Settlor}: Person who creates trust
				\item \textbf{Trustee}: Person/entity managing trust
				\item \textbf{Protector}: If has veto or replacement rights
				\item \textbf{Beneficiaries}: To extent determinable
			\end{itemize}
		\end{block}
		\textcolor{red}{Protector with veto qualifies as PEP if former head of state.}
	\end{frame}
	
	\begin{frame}{A.14.2: Bearer Shares - Highest Vulnerability}
		\begin{block}{Bearer Shares Risk}
			\begin{itemize}
				\item Transferred by physical delivery of certificate
				\item \textbf{NO public registry record} of ownership
				\item Ownership can be transferred anonymously
				\item FATF Rec. 24: prohibit or require immobilization
			\end{itemize}
		\end{block}
		\textbf{Bank obligations:}
		\begin{itemize}
			\item Verify custodian can identify holder
			\item Still need to identify UBO
			\item Decline account if custodian refuses disclosure
		\end{itemize}
	\end{frame}
	
	% ============================================================
	% SECTION A.15: SAR FILING AND TIPPING-OFF
	% ============================================================
	
	\begin{frame}{A.15.1: SAR Filing Standard - Reasonable Suspicion}
		\begin{block}{When to File a SAR}
			\begin{itemize}
				\item Standard = \textbf{reasonable suspicion}
				\item \textbf{No requirement for conclusive proof}
				\item File \textbf{immediately upon suspicion}
				\item Required regardless of transaction value
				\item Failure to file = regulatory risk
				\item BSA Officer can face personal liability
			\end{itemize}
		\end{block}
		\textcolor{red}{"Evidence not conclusive" is NOT a valid reason to decline filing.}
	\end{frame}
	
	\begin{frame}{A.15.2: Tipping-Off - Prohibited vs. Permissible Disclosures}
		\begin{block}{Tipping-Off Prohibitions}
			\textbf{Prohibited (criminal penalties):}
			\begin{itemize}
				\item Telling customer a SAR was filed
				\item Disclosing compliance is "reviewing transactions"
				\item Disclosing investigation to third parties
				\item Publicly announcing customer under investigation
			\end{itemize}
			\textbf{Permissible:}
			\begin{itemize}
				\item Internal sharing with legal department
				\item 314(b) sharing with other opted-in institutions
				\item Responding to law enforcement subpoena
			\end{itemize}
		\end{block}
	\end{frame}
	
	\begin{frame}{A.15.3: Law Enforcement "Keep Open" Requests}
		\begin{block}{Protocol for Keep Open Requests}
			\begin{enumerate}
				\item Obtain request \textbf{IN WRITING} from law enforcement
				\item Document thoroughly
				\item Obtain senior management approval
				\item Maintain heightened monitoring
				\item Do NOT disclose to customer
			\end{enumerate}
		\end{block}
		\textbf{Do NOT:} Publicly announce, ignore request, close account without coordination.
	\end{frame}
	
	% ============================================================
	% SECTION A.16: NPO TERRORIST FINANCING VULNERABILITIES
	% ============================================================
	
	\begin{frame}{A.16.1: NPO Terrorist Financing Red Flags}
		\begin{block}{FATF Recommendation 8 - NPO Vulnerabilities}
			\textbf{Vulnerabilities:}
			\begin{itemize}
				\item Operate in conflict zones with weak monitoring
				\item Receive large cash donations with limited scrutiny
				\item Can receive funds from Grey/Black List jurisdictions
			\end{itemize}
			\textbf{Red flags requiring SAR:}
			\begin{itemize}
				\item Large donation from Grey List jurisdiction
				\item Rapid cash withdrawal (80\%+ within 48 hours)
				\item Cash withdrawal far from headquarters
				\item Transfers to conflict zone sub-NPOs where terrorists control aid
				\item No documentation matching cash withdrawals to mission
			\end{itemize}
		\end{block}
	\end{frame}
	
% ============================================================
% SECTION A.17: HIGH-DIFFICULTY MIXED SCENARIOS
% ============================================================

\begin{frame}{A.17.1: Layered Sovereign Corruption Extraction}
	\begin{block}{Sovereign Corruption Scheme Analysis}
		\textbf{Scheme:} State-owned oil company signs "service agreements" with BVI shell company. BVI-Co receives \$220M. Funds split: \$150M to Cayman trust (PEP's son beneficiary), \$70M to London real estate.
	\end{block}
	\textbf{Obligations:}
	\begin{itemize}
		\item Immediate escalation to senior management
		\item Emergency EDD review
		\item Assess if transactions constitute corruption proceeds
		\item File STR if reasonable suspicion exists
		\item Cannot close account without approval
		\item Ensure no tipping-off
	\end{itemize}
\end{frame}

\begin{frame}{A.17.2: Cross-Border Cash Courier Controls}
	\begin{block}{FATF Recommendation 32 - Cash Courier}
		\textbf{Deficiencies:}
		\begin{itemize}
			\item No declaration system for cross-border currency
			\item Cash couriers intercepted with \$500,000+ but no penalties
		\end{itemize}
		\textbf{Required controls:}
		\begin{itemize}
			\item Declaration OR disclosure system at borders
			\item Authority to stop, restrain, question couriers
			\item Sanctions for false declarations
		\end{itemize}
	\end{block}
\end{frame}

\begin{frame}{A.17.3: Money Mule Network (P2P Platform)}
	\begin{block}{P2P Lending Platform Exploitation}
		\textbf{Pattern:}
		\begin{itemize}
			\item 17 accounts registered within 48-hour window, same VPN
			\item Each received small loans from single institutional lender
			\item All repaid within 3 days using offshore wires
			\item After repayment, withdrew to prepaid debit cards
			\item Institutional lender is anonymous LLC
		\end{itemize}
		\textbf{Typology:} Self-funding layering ring.
	\end{block}
\end{frame}

\begin{frame}{A.17.4: FATF Recommendation 24 - Legal Person Transparency}
	\begin{block}{Beneficial Ownership Registry Requirements}
		\textbf{Deficiencies:}
		\begin{itemize}
			\item BO registry not accessible to FIs without court order
			\item Law enforcement access requires court order
			\item No requirement to update BO information
		\end{itemize}
		\textbf{Required:}
		\begin{itemize}
			\item Adequate, accurate, CURRENT BO information
			\item Accessible to competent authorities in timely manner
			\item Mechanisms to ensure accuracy and updates
		\end{itemize}
	\end{block}
\end{frame}

\begin{frame}{A.17.5: FATF Recommendation 30 - Financial Investigations}
	\begin{block}{Financial Investigation Deficiencies}
		\textbf{Problem:}
		\begin{itemize}
			\item LE cannot obtain financial records directly from banks
			\item All requests go through prosecutor (6-12 months)
			\item Assets moved before records obtained
		\end{itemize}
		\textbf{Required:}
		\begin{itemize}
			\item Power to obtain records expeditiously (days, not months)
			\item Ability to identify, trace, freeze assets
			\item Parallel financial investigations
		\end{itemize}
	\end{block}
\end{frame}

\begin{frame}{A.17.6: FATF Immediate Outcome 5 - Preventive Measures}
	\begin{block}{IO.5 Assessment}
		\textbf{Deficiencies:}
		\begin{itemize}
			\item No beneficial ownership registry
			\item No requirement for companies to identify UBOs
			\item No requirement to update BO information
		\end{itemize}
		\textbf{Impact:}
		\begin{itemize}
			\item FIs cannot verify UBO information
			\item Undermines preventive measures
			\item Information quickly outdated
		\end{itemize}
	\end{block}
\end{frame}

\begin{frame}{A.17.7: FATF Immediate Outcome 2 - International Cooperation}
	\begin{block}{IO.2 Assessment}
		\textbf{Deficiencies:}
		\begin{itemize}
			\item Never submitted or responded to MLAT request
			\item No designated central authority for MLAT processing
		\end{itemize}
		\textbf{Why deficient:}
		\begin{itemize}
			\item MLATs are primary mechanism for obtaining evidence
			\item Egmont exchanges CANNOT replace MLATs
			\item Absence of central authority creates delays
		\end{itemize}
	\end{block}
\end{frame}

\begin{frame}{A.17.8: FATF Recommendation 34 - Guidance and Feedback}
	\begin{block}{Recommendation 34}
		\textbf{Deficiency:} No AML/CFT guidance issued to FIs; no feedback on STR quality.
		\vspace{0.2cm}
		\textbf{Required:}
		\begin{itemize}
			\item Guidance and feedback to assist FIs
			\item Red flag indicators, typologies, emerging trends
			\item Feedback on STR quality
		\end{itemize}
		\textbf{This is NOT optional} - Rec. 34 is a standard, not a best practice.
	\end{block}
\end{frame}

\begin{frame}{A.17.9: FATF Recommendation 14 - MVTS / Hawala}
	\begin{block}{Money or Value Transfer Services Regulation}
		\textbf{Deficiency:} Unlicensed hawala operators; government chose not to regulate.
		\vspace{0.2cm}
		\textbf{Required (Rec. 14):}
		\begin{itemize}
			\item All MVTS providers MUST be licensed or registered
			\item Subject to AML/CFT monitoring
			\item Unlicensed providers subject to sanctions
			\item \textbf{No cultural heritage exemption} for hawala
		\end{itemize}
	\end{block}
\end{frame}

\begin{frame}{A.17.10: FATF Recommendation 26 - Supervision of FIs}
	\begin{block}{Supervisory Deficiencies}
		\textbf{Findings:}
		\begin{itemize}
			\item On-site exams every 5 years for ALL banks (same frequency)
			\item No off-site monitoring or thematic reviews
			\item No penalties issued in 10 years
		\end{itemize}
		\textbf{Required:}
		\begin{itemize}
			\item Risk-based approach to exam frequency
			\item On-site AND off-site monitoring
			\item Sanctions must be imposed for violations
		\end{itemize}
	\end{block}
\end{frame}

% ============================================================
% CONCLUSION SLIDES
% ============================================================

\begin{frame}{Domain A Summary: Top 20 Most Tested Concepts}
	\begin{enumerate}
		\item \textbf{ML vs. TF distinction}
		\item \textbf{Three stages}: Placement, Layering, Integration
		\item \textbf{TBML techniques}: Over/under-invoicing, phantom shipping
		\item \textbf{BMPE}: Uses drug dollars to pay legitimate invoices
		\item \textbf{Nested correspondent accounts}: Loss of visibility
		\item \textbf{Shell banks}: No physical presence; prohibited
		\item \textbf{PEPs}: Family members included; SoW AND SoF required
		\item \textbf{6AMLD}: 22 predicate offenses; corporate liability
		\item \textbf{Travel Rule for VASPs}: Receiving VASP has obligations
		\item \textbf{DeFi}: Control or influence = VASP obligations
		\item \textbf{Privacy coins}: Ring signatures, stealth addresses
		\item \textbf{NFT wash trading}: Circular self-dealing
		\item \textbf{Human trafficking}: Multiple individuals to one controller
		\item \textbf{Chip washing}: Large buy-in, minimal gaming
		\item \textbf{IOLTA vulnerabilities}: Banks don't look through
		\item \textbf{Bearer shares}: Transfer by physical delivery
		\item \textbf{SAR filing standard}: Reasonable suspicion
		\item \textbf{Tipping-off}: Telling customer about SAR
		\item \textbf{De-risking}: Blanket termination violates RBA
		\item \textbf{FATF effectiveness}: IO.6, IO.7, IO.8
	\end{enumerate}
\end{frame}

\begin{frame}{Critical Red Flags - Quick Reference}
	\begin{itemize}
		\item \textbf{Structuring:} \$9,000-\$9,800 repeated deposits
		\item \textbf{PEP:} Minister's daughter, "consulting fees," refuses SoW docs
		\item \textbf{TBML:} Price 400\% above/below market, shell company exporter
		\item \textbf{Human Trafficking:} Multiple individuals $\rightarrow$ 1 controller $\rightarrow$ multiple utility payments
		\item \textbf{NPO/TF:} Rapid cash withdrawal far from HQ
		\item \textbf{Correspondent:} Undisclosed nested transactions
		\item \textbf{DeFi:} 68\% governance token control, ability to upgrade code
		\item \textbf{Sanctions:} Fuzzy name matching required
		\item \textbf{Casino:} Large cash buy-in, minimal gaming, cash-out check
		\item \textbf{Real Estate:} Recently formed LLC, IOLTA wire, refuses UBO disclosure
	\end{itemize}
\end{frame}

\begin{frame}{Exam Day Reminders for Domain A}
	\begin{itemize}
		\item \textbf{Domain A is 30\% of the exam} - largest single domain
		\item FATF Recommendations, PEPs, TBML, correspondent banking heavily tested
		\item \textbf{Select Two/Three/Four questions} = get ALL correct for credit
		\item Read scenarios carefully - facts matter
		\item Distinguish between ML and TF
		\item Know placement, layering, integration stages
		\item SARs filed on reasonable suspicion, not conclusive proof
		\item De-risking (blanket termination) is NOT preferred
		\item OSINT/adverse media does NOT require conviction
	\end{itemize}
	\vspace{0.5cm}
	\centering
	\textcolor{blue}{\Large Good luck on your CAMS exam!}
\end{frame}

\end{document}